\documentclass[11pt,a4paper]{article}

% ── Packages (only those available on this installation) ──────────────────────
\usepackage[margin=2.5cm]{geometry}
\usepackage{amsmath,amssymb}
\usepackage{graphicx}
\usepackage{booktabs}
\usepackage{longtable}
\usepackage{array}
\usepackage{xcolor}
\usepackage{hyperref}
\usepackage{caption}
\usepackage{subcaption}
\usepackage{enumitem}
\usepackage{microtype}
\usepackage{listings}
\usepackage{float}
\usepackage{placeins}

% ── Figure path ───────────────────────────────────────────────────────────────
\graphicspath{{figures/}}

% ── Hyperref setup ────────────────────────────────────────────────────────────
\hypersetup{
  colorlinks = true,
  linkcolor  = blue!70!black,
  citecolor  = green!50!black,
  urlcolor   = blue!70!black,
}

% ── Custom commands ───────────────────────────────────────────────────────────
\newcommand{\numu}{$\nu_\mu$}
\newcommand{\numubar}{$\bar{\nu}_\mu$}
\newcommand{\ccpip}{$\mathrm{CC}\pi^{\pm}$}
\newcommand{\pmu}{$p_\mu$}
\newcommand{\ppi}{$p_\pi$}
\newcommand{\costhmu}{$\cos\theta_\mu$}
\newcommand{\costhpi}{$\cos\theta_\pi$}
\newcommand{\thmupi}{$\theta_{\mu\pi}$}
\newcommand{\uboone}{MicroBooNE}
\newcommand{\GeVc}{\,\text{GeV}\!/c}
\newcommand{\pot}{\,\text{POT}}
\newcommand{\cm}{\,\text{cm}}

% inline math shorthand so they work inside math mode too
\newcommand{\pmuM}{p_\mu}
\newcommand{\ppiM}{p_\pi}
\newcommand{\cthmuM}{\cos\theta_\mu}
\newcommand{\cthpiM}{\cos\theta_\pi}
\newcommand{\thmupM}{\theta_{\mu\pi}}

% ── Listings (C++ snippets) ───────────────────────────────────────────────────
\lstset{
  language=C++,
  basicstyle=\ttfamily\small,
  keywordstyle=\bfseries\color{blue!70!black},
  commentstyle=\itshape\color{gray},
  backgroundcolor=\color{gray!10},
  frame=single,
  framerule=0.4pt,
  numbers=none,
  breaklines=true,
}

% ── Table column helpers ──────────────────────────────────────────────────────
\newcolumntype{C}[1]{>{\centering\arraybackslash}p{#1}}
\newcolumntype{R}[1]{>{\raggedleft\arraybackslash}p{#1}}
\newcolumntype{L}[1]{>{\raggedright\arraybackslash}p{#1}}

% Let TeX stretch spacing on hard-to-break lines (long \texttt paths/code) to keep
% text out of the margin.
\emergencystretch=3em

% ─────────────────────────────────────────────────────────────────────────────

\usepackage{xr}
\externaldocument[N-]{analysis_note}
\externaldocument[CR-]{cr_data_note}

\newcommand{\GeVcc}{\,\text{GeV}\!/c}

\begin{document}

\begin{titlepage}
\centering
\vspace*{2cm}
{\LARGE\bfseries Binning criterion, proton-tagged angular observables\\ and two-dimensional cross sections\par}
\vspace{0.5cm}
{\large Measurement of $\nu_\mu + \bar{\nu}_\mu$ Charged-Current Single Pion
Production on Argon in the NuMI Beam\par}
\vspace{0.8cm}
{\large Studies of 2026-09-16/17; blind (fake data throughout)\par}
\vspace{0.5cm}
{\normalsize 2026-09-17\par}
\vspace{1.5cm}
{\large Ishanee Pophale, Jarek Nowak\par}
\vspace{0.5cm}
{\normalsize\textit{Lancaster University}\par}
\vspace{2cm}

\begin{abstract}
Three questions were asked of the released analysis: whether the migration-diagonal requirement
could be relaxed from $0.68$ to $0.50$, whether the proton-tagged sample should measure the proton
angle and the pion--proton opening angle, and whether double-differential cross sections are
supported. This note records the answers and the machinery built for them. \textbf{(i)}~Relaxing the
criterion admits finer binnings in $\cos\theta_\mu$, $\theta_\mu$, $p_\pi$ and the proton-tagged
transverse variables, but every one of the twenty-six candidates built and unfolded in full is
rejected by the conditioning of the additional-smearing matrix $A_C$: the criterion is not what
limits this measurement, and the study is reported as a study, with no rebinning adopted.
\textbf{(ii)}~$\theta_p$ (four bins) and the pion--proton opening angle $\theta_{\pi p}$ (three bins)
both meet the unchanged $0.68$ criterion and are now extracted for FHC, RHC and the combined
configuration, with generator predictions from all four generators; $\theta_{\pi p}$ is the
best-conditioned new observable of the analysis ($s_{\min}/s_1 = 0.30$--$0.48$). The selection now
stores the opening angle, which required reprocessing the proton-tagged tree; the reprocess was
gated on a bit-for-bit comparison of every existing branch. \textbf{(iii)}~The framework supports
two-dimensional measurements as flattened bins with one slice per outer bin; eight pairs were
screened and are being extracted. The first results confirm what the one-dimensional study
already showed: a good migration diagonal does not guarantee a usable $A_C$. The joint FHC$+$RHC
background fit developed in the same period is documented separately in the control-region note.
\end{abstract}
\end{titlepage}

\tableofcontents
\clearpage

% =============================================================================
\section{Scope, status and conventions}
\label{sec:scope}
% =============================================================================

The signal region is blind: every ``measurement'' here is the per-run Poisson fake data of the
released throw, and the extractions are closure tests. Three families of configuration are used and
must not be confused.
\begin{itemize}\itemsep2pt
  \item \textbf{Live}: $\theta_p$ and $\theta_{\pi p}$, in \texttt{configs/ccpi1p\_*} with the released
    naming, extracted into \texttt{processed/} beside the released observables. They use binnings that
    meet the unchanged $0.68$ criterion.
  \item \textbf{0.50 study}: \texttt{configs/c50/}, extracted into \texttt{processed/rebuild\_c50/}.
    Nothing from it is adopted. The directory is deliberately outside the glob of
    \texttt{norm\_manifest.py}, which audits the release.
  \item \textbf{2D study}: \texttt{configs/d2/}, extracted into \texttt{processed/rebuild\_2d/}.
\end{itemize}
Screening uses the per-run exposure scales and the safe central-value weight of the extraction, and
the same first/last-edge conventions as the released bin configs. Two criteria appear throughout: the
column-normalised migration diagonal per true bin (the BNB criterion, $>0.68$) and the conditioning
of the additional-smearing matrix, $s_{\min}/s_1$ from its singular values, which the analysis note
already identified as the quantity that decided $\cos\theta_\pi$ and $\theta_{\mu\pi}$.

% =============================================================================
\section{The 0.50 migration-diagonal study}
\label{sec:c50}
% =============================================================================

\subsection{Screening}

\texttt{macros/binning\_screen\_2d.C} evaluates every candidate in one pass over the processed files
(the earlier \texttt{binning\_screen.C} used four \texttt{TTree::Draw} calls per bin, which does not
scale to a long candidate list). It reproduces the published pre-screen table exactly for FHC and RHC,
which is its regression test, and it also handles two-dimensional candidates
(\S\ref{sec:2d}). For the proton-tagged observables the maximin edge scan of
\texttt{macros/tki\_binning\_scan.C} was generalised to take the criterion, an observable filter and
the maximum bin count, and to cover $W$, the proton angles and the opening angle; it reproduces the
two-bin transverse-kinematic-imbalance numbers of the binning README.

At $0.50$ the following become admissible where $0.68$ refused them (smallest diagonal, FHC / RHC):
$\cos\theta_\mu$ with 6, 7 or 8 bins ($0.64/0.63$, $0.61/0.61$, $0.60/0.59$); $\theta_\mu$ with 7 or 8
($0.63/0.62$, $0.61/0.61$); $p_\pi$ with three bins and the extra edge at $0.26$\GeVcc\ ($0.51/0.53$);
proton-tagged $\delta p_T$ with three ($0.59/0.55$), $\delta\alpha_T$ with three ($0.52/0.50$),
$W_{\pi p}$ with two at $1.19$\GeVcc\ ($0.56/0.56$); $\cos\theta_p$ with six and $\theta_p$ with seven.
$p_\mu$ still fails ($0.33/0.31$ released, worse when finer), as do three-bin $p_\pi$ at $0.30$\GeVcc,
four-bin $p_n$ and $\delta\phi_T$ in RHC. $W_{\rm had}$ in two bins passes even at $0.68$
($0.71/0.72$) although the released configuration has six.

\subsection{What the full extractions say}

Twenty-six candidates (thirteen binnings $\times$ FHC and RHC) were built and unfolded with the
released machinery. Table~\ref{tab:c50} gives the result. The conditioning is decisive: of the finer
binnings only $W_{\rm had}$ (two bins) keeps an $A_C$ with $s_{\min}/s_1$ above $0.5$, and the
two-bin $W_{\pi p}$ and the four-bin $\cos\theta_p$ in FHC are the only others above $0.1$. Several
candidates have $s_{\min}/s_1$ numerically zero to three decimals: their $A_C$ maps every model onto
a single shape, so the extraction cannot be read as a shape measurement, exactly the pathology the
analysis note documented for $\theta_\mu$. Closure $\chi^2$ values are uniformly small --- which is
the point: they cannot distinguish these cases, because a rank-deficient $A_C$ makes data and
prediction agree by construction.

\begin{table}[H]
  \centering
  \caption{Candidates admitted by the $0.50$ criterion, built and unfolded in full (fake data).
    $A_C$ rows: range of the row sums; $s_{\min}/s_1$: conditioning of $A_C$; data stat.: bin-averaged
    data-statistical term; unc.: range of the per-bin total uncertainty; closure: range of the
    per-bin ratio to the $A_C$-smeared truth. Rows in bold keep a usable $A_C$.}
  \label{tab:c50}
  \scriptsize\setlength{\tabcolsep}{3pt}
  \begin{tabular}{llccccc@{\hspace{4pt}}c}
    \toprule
    Candidate & Config & bins & $A_C$ rows & $s_{\min}/s_1$ & data stat. & unc.\ [\%] & closure \\
    \midrule
    $\cos\theta_\mu$ 6 & FHC & 6 & $0.40$--$0.52$ & $0.110$ & $17.0\%$ & $30$--$76$ & $0.80$--$1.21$ \\
                       & RHC & 6 & $0.09$--$0.44$ & $0.085$ & $15.8\%$ & $31$--$133$ & $0.58$--$1.25$ \\
    $\cos\theta_\mu$ 7 & FHC & 7 & $0.21$--$0.37$ & $0.024$ & $23.8\%$ & $31$--$143$ & $0.51$--$1.30$ \\
                       & RHC & 7 & $0.03$--$0.52$ & $0.054$ & $19.3\%$ & $33$--$121$ & $0.71$--$1.17$ \\
    $\theta_\mu$ 7 & FHC & 7 & $0.47$--$0.68$ & $0.000$ & $11.0\%$ & $25$--$52$ & $1.02$--$1.38$ \\
                   & RHC & 7 & $0.35$--$0.93$ & $0.001$ & $10.3\%$ & $30$--$38$ & $1.03$--$1.38$ \\
    $\theta_\mu$ 8 & FHC & 8 & $0.53$--$0.73$ & $0.000$ & $10.0\%$ & $25$--$54$ & $0.87$--$1.36$ \\
                   & RHC & 8 & $0.16$--$0.76$ & $0.001$ & $13.2\%$ & $28$--$54$ & $0.84$--$1.42$ \\
    $p_\pi$ 3 & FHC & 3 & $0.53$--$0.91$ & $0.048$ & $14.8\%$ & $39$--$56$ & $0.90$--$0.97$ \\
              & RHC & 3 & $0.45$--$0.89$ & $0.010$ & $12.7\%$ & $40$--$45$ & $0.97$--$1.02$ \\
    $\delta p_T$ 3 & FHC & 3 & $0.37$--$0.89$ & $0.001$ & $13.8\%$ & $32$--$37$ & $0.94$--$0.96$ \\
                   & RHC & 3 & $0.32$--$0.86$ & $0.000$ & $12.0\%$ & $35$--$36$ & $1.04$--$1.12$ \\
    $\delta\alpha_T$ 3 & FHC & 3 & $0.44$--$0.98$ & $0.029$ & $10.2\%$ & $42$--$52$ & $0.83$--$0.99$ \\
                       & RHC & 3 & $0.42$--$0.92$ & $0.011$ & $9.8\%$ & $41$--$47$ & $1.00$--$1.17$ \\
    $W_{\pi p}$ 2 ($1.19$) & FHC & 2 & $0.27$--$0.91$ & $0.264$ & $19.1\%$ & $41$--$72$ & $0.93$--$0.99$ \\
                           & RHC & 2 & $0.28$--$1.48$ & $0.147$ & $24.1\%$ & $47$--$91$ & $0.73$--$1.00$ \\
    \textbf{$W_{\rm had}$ 2} & FHC & 2 & $0.80$--$0.90$ & $\mathbf{0.579}$ & $16.3\%$ & $37$--$87$ & $0.80$--$1.03$ \\
                             & RHC & 2 & $0.76$--$0.86$ & $\mathbf{0.541}$ & $17.1\%$ & $45$--$100$ & $0.78$--$1.05$ \\
    $\cos\theta_p$ 4 & FHC & 4 & $0.54$--$0.76$ & $0.445$ & $20.6\%$ & $34$--$76$ & $0.66$--$1.08$ \\
                     & RHC & 4 & $-0.09$--$0.62$ & $0.018$ & $36.2\%$ & $41$--$182$ & $0.43$--$1.41$ \\
    $\cos\theta_p$ 6 & FHC & 6 & $0.77$--$0.90$ & $0.000$ & $25.6\%$ & $34$--$84$ & $0.54$--$1.06$ \\
                     & RHC & 6 & $0.43$--$0.78$ & $0.000$ & $24.3\%$ & $38$--$72$ & $0.75$--$1.48$ \\
    $\theta_p$ 5 & FHC & 5 & $0.59$--$1.14$ & $0.034$ & $13.5\%$ & $33$--$44$ & $0.84$--$1.24$ \\
                 & RHC & 5 & $-0.03$--$0.60$ & $0.005$ & $24.9\%$ & $39$--$99$ & $0.82$--$1.30$ \\
    $\theta_p$ 7 & FHC & 7 & $0.61$--$0.98$ & $0.000$ & $16.4\%$ & $34$--$58$ & $0.87$--$1.30$ \\
                 & RHC & 7 & $0.32$--$0.87$ & $0.000$ & $20.7\%$ & $39$--$60$ & $0.77$--$1.14$ \\
    \bottomrule
  \end{tabular}
\end{table}

\paragraph{Conclusion.} The $0.68$ criterion is not the binding constraint of this analysis and
lowering it to $0.50$ buys nothing: the binnings it admits are rejected one level down, by the
regularised response. No rebinning is proposed. Two by-products are worth keeping: $W_{\rm had}$ in
two bins is better conditioned than the released six-bin configuration, and the criterion study
produced the screening tool used for the new observables below.

\begin{figure}[H]
  \centering
  \includegraphics[width=0.92\textwidth]{c50/diag_incl_fhc}
  \caption{Migration diagonal per true bin for the released inclusive binnings and the finer
    candidates, FHC; the two criteria are drawn. The RHC version and the proton-tagged and
    two-dimensional versions are \texttt{c50/diag\_\{incl,1p,2d\}\_\{fhc,rhc\}}.}
  \label{fig:diag_incl}
\end{figure}

% =============================================================================
\section{Proton-tagged angular observables}
\label{sec:newobs}
% =============================================================================

\subsection{What was added, and why the opening angle needed a reprocess}

The proton-tagged selection already stored the proton polar angle about the neutrino direction
($\cos\theta_p$, reconstructed in the beta convention, truth against the per-event neutrino). It did
not store the pion--proton opening angle $\theta_{\pi p}$, nor the index of the chosen proton, so the
angle could not be formed offline from the stored quantities. Both are now written by
\texttt{CC1mu1pi1p}: the reconstructed angle from the two track directions (frame independent, and
immune to the zero-momentum degeneracy of a scaled vector), the true angle from the leading charged
pion and the leading proton.

Before committing to the reprocess, the proton choice was replicated offline
(\texttt{macros/proton\_\allowbreak angle\_\allowbreak friend.C}) --- the lowest-LLR track above threshold, contained, not the
muon or pion candidate --- and validated against the stored $\cos\theta_p$: identical for all
$18\,698$ selected events of the fourteen Monte Carlo files. That friend tree also answered what the
candidate is: on selected \emph{signal} events it is the leading true proton $82\%$ of the time, a
softer proton $6\%$, and not a proton $13\%$ (on all selected events, background included, $59$,
$10$ and $31\%$). The wrong-candidate events populate a flat tail in both angles, while the correct
ones peak within about $0.1$~rad (figure \texttt{c50/resolution\_\allowbreak proton\_\allowbreak angles}).

\paragraph{Reprocess and its gate.} The proton-tagged tree (41 units) and the seven per-run beam-off
files, which are shared with the inclusive tree and carry both selections, were reprocessed into
staging directories, and the per-run fake data re-thrown with the same seeds. Promotion was gated on
\texttt{macros/compare\_reprocess.C}: for every file, all 168 numeric branches of the old tree
identical entry by entry, the entry count and \texttt{summed\_pot} unchanged, and the new angle equal
to the friend-tree value on every selected and signal event. All 56 files passed. One needed
attention: the FHC Run-1 dirt snapshot carries a \texttt{summed\_pot} written in place by
\texttt{set\_summed\_pot.C}, $1.673920\times10^{21}$ against the native $1.6739168\times10^{21}$; the
established value was carried into the new file so that the normalisation cache keys are unchanged.
The old trees are kept as \texttt{w\_pre\_thpipr} and \texttt{ext\_perrun\_pre\_thpipr}.

\subsection{Binning}

The maximin scan (FHC Runs 1, 2, 4; $\geq400$ events per bin) gives the largest bin count at each
criterion; the weighted screen then checks it over FHC and RHC at the released exposure
(Table~\ref{tab:newbins}). $\theta_p$ supports one bin more than $\cos\theta_p$ at both criteria, so
the angle rather than its cosine is measured, matching the released $\theta_\mu$ convention. Both
adopted binnings meet $0.68$; the RHC $\theta_p$ value sits exactly at it, as the released
$\cos\theta_\mu$ does.

\begin{table}[H]
  \centering
  \caption{Proton-tagged angular binnings. Scan: worst diagonal of the maximin scheme at that bin
    count (unweighted, FHC Runs 1, 2, 4). Screen: smallest diagonal FHC / RHC with the per-run
    weighting. Adopted binnings in bold.}
  \label{tab:newbins}
  \small
  \begin{tabular}{llccl}
    \toprule
    Observable & bins & scan & screen FHC / RHC & edges \\
    \midrule
    \textbf{$\theta_p$} & \textbf{4} & $\mathbf{72.6\%}$ & $\mathbf{0.73/0.68}$ & $0$, $0.597$, $0.942$, $1.287$, $\pi$ \\
    $\theta_p$ & 5 & $68.0\%$ & $0.67/0.64$ & $0$, $0.377$, $0.691$, $0.942$, $1.225$, $\pi$ \\
    $\theta_p$ & 7 & $57.1\%$ & $0.56/0.55$ & (0.50 criterion) \\
    $\cos\theta_p$ & 4 & $73.2\%$ & $0.73/0.66$ & $-1$, $0.275$, $0.575$, $0.825$, $1$ \\
    $\cos\theta_p$ & 6 & $57.2\%$ & $0.58/0.55$ & (0.50 criterion) \\
    \textbf{$\theta_{\pi p}$} & \textbf{3} & $\mathbf{75.8\%}$ & $\mathbf{0.77/0.75}$ & $0$, $1.193$, $2.010$, $\pi$ \\
    $\theta_{\pi p}$ & 4 & $67.9\%$ & --- & $0$, $0.973$, $1.476$, $2.010$, $\pi$ \\
    $\theta_{\pi p}$ & 7 & $53.9\%$ & $0.52/0.53$ & (0.50 criterion) \\
    \bottomrule
  \end{tabular}
\end{table}

\subsection{Generator predictions}

The four generator predictions are produced in \texttt{generator\_predictions/gen2d/}, which holds
copies of the released readers (GENIE gst, GiBUU \texttt{FinalEvents}, NuWro and NEUT in their
containers) with the new observables added through one shared header, so the released chain is
untouched. Two checks are built in. Every histogram the copies also write is reproduced bit-for-bit
from the released per-flavour files, for all four generators; and the generic combiner reproduces all
sixty released inclusive FTE histograms (four generators $\times$ FHC, RHC, combined $\times$ five
observables) exactly. Unlike the released proton-tagged predictions, which exist for FHC only, the new
ones are produced for FHC, RHC and the combined configuration.
\texttt{scripts/check\_ext\_bins.py} checks that the generator edges and the configuration edges agree
and that each prediction has as many bins as its configuration.

\subsection{Extractions}

\begin{table}[H]
  \centering
  \caption{Extractions of the new observables (fake data). Columns as in Table~\ref{tab:c50};
    $\chi^2$ is the closure against the $A_C$-smeared truth.}
  \label{tab:newobs_extract}
  \small
  \begin{tabular}{llcccccc}
    \toprule
    Observable & Config & bins & $A_C$ rows & $s_{\min}/s_1$ & data stat. & unc.\ [\%] & closure $\chi^2$/ndf \\
    \midrule
    $\theta_p$ & FHC & 4 & $0.58$--$1.30$ & $0.082$ & $12.3\%$ & $35$--$45$ & $0.91/4$ \\
               & RHC & 4 & $0.21$--$1.00$ & $0.002$ & $12.9\%$ & $36$--$44$ & $0.49/4$ \\
               & COMB & 4 & $0.60$--$1.24$ & $0.125$ & $9.7\%$ & $32$--$50$ & $0.76/4$ \\
    $\theta_{\pi p}$ & FHC & 3 & $0.57$--$0.65$ & $0.480$ & $20.1\%$ & $38$--$93$ & $0.16/3$ \\
                     & RHC & 3 & $0.59$--$0.73$ & $0.376$ & $24.5\%$ & $43$--$187$ & $0.12/3$ \\
                     & COMB & 3 & $0.43$--$0.49$ & $0.304$ & $14.3\%$ & $40$--$89$ & $0.06/3$ \\
    \bottomrule
  \end{tabular}
\end{table}

$\theta_{\pi p}$ is the best-conditioned observable of the proton-tagged programme: its
$s_{\min}/s_1$ of $0.30$--$0.48$ is above everything in Table~\ref{tab:c50} except two-bin
$W_{\rm had}$, and above the released $\cos\theta_\pi$ ($0.27/0.24$). Its price is the last bin, whose
total uncertainty reaches $93\%$ (FHC) and $187\%$ (RHC). $\theta_p$ unfolds well in FHC and the
combined configuration but is nearly rank-deficient in RHC ($0.002$), where its diagonal sat exactly
at the criterion: the RHC $\theta_p$ curves should not be read as a shape measurement.

\begin{figure}[H]
  \centering
  \includegraphics[width=\textwidth]{bkgfit/newobs_xsec}
  \caption{Extracted $\mathrm{d}\sigma/\mathrm{d}\theta_p$ and
    $\mathrm{d}\sigma/\mathrm{d}\theta_{\pi p}$ (fake data; the signal region is blind) against the
    four generators and the fake-data truth. Integrated cross sections in the legends,
    $10^{-38}$~cm$^2$/Ar, $A_C$-smeared.}
  \label{fig:newobs_xsec}
\end{figure}

Both distributions are strongly forward peaked: the top $\theta_p$ bin, more than half the angular
range, carries about a tenth of the first bin's $\mathrm{d}\sigma/\mathrm{d}\theta$, and the pion and
proton are separated by less than $2$~rad in the bulk of the events. The generator ordering is the
familiar one --- GENIE lowest, NEUT highest, GiBUU and NuWro between, a spread of about $1.6$ in the
integral --- and the shapes are similar, so these observables discriminate mainly through
normalisation, with NEUT's relative excess in the middle $\theta_{\pi p}$ bin the exception.

% =============================================================================
\section{Two-dimensional cross sections}
\label{sec:2d}
% =============================================================================

\subsection{What the framework already supports}

A double-differential measurement needs no new unfolding machinery. UniverseMaker bins are arbitrary
selection strings, so a two-dimensional binning is a flat list of bins; \texttt{SliceBinning} carries
slices with an ``other variable'' held in a fixed range, and \texttt{UnfolderNuMI} loops over the
slices and divides by the width of the other variable, so each slice is a differential cross section
in the inner variable. Two constraints were found and respected: the analysis bins of a slice must be
contiguous, because the per-slice $A_C$ block is taken as a range, so the bin configuration is written
slice by slice; and the closure sidecar was dumped only for slice~0, which for a two-dimensional
configuration is the first outer bin rather than the whole measurement, so \texttt{UnfolderNuMI} now
also dumps the ``bin number'' slice as \texttt{closure\_hists\_all\_*} when a configuration has more
than one physical slice (one-dimensional output is unchanged). Configurations are generated by
\texttt{scripts/make\_bin\_configs.py}, which reproduces the hand-written one-dimensional candidate
configurations byte for byte as its regression test.

\subsection{Screening}

Eight pairs were screened with \texttt{binning\_screen\_2d.C}, which for a two-dimensional candidate
also splits the off-diagonal migration into leakage within the outer slice (X) and between slices (Y).
Table~\ref{tab:2dscreen} gives the best grid of each pair. The Y-leakage is what decides: the pairs
built on $\delta p_T$ lose $20$--$30\%$ of events across the outer bin and those built on $W_{\pi p}$
up to $60\%$, which is why every $W_{\pi p}$ pair fails both criteria. Expected event counts are the
second constraint: $\cos\theta_\pi\times\theta_{\mu\pi}$ has a bin with four expected events.

\begin{table}[H]
  \centering
  \caption{Two-dimensional pairs (inner $\times$ outer), best grid of each: smallest diagonal over the
    flattened true bins, mean leakage within and between outer slices, and the expected selected
    signal of the thinnest bin at data exposure.}
  \label{tab:2dscreen}
  \small
  \begin{tabular}{llcccc}
    \toprule
    Pair & grid & min.\ diagonal FHC / RHC & X-leak & Y-leak & thinnest bin \\
    \midrule
    $\cos\theta_\pi\times\cos\theta_\mu$ & $2\times3$ & $0.69/0.68$ & $3.4\%$ & $12.5\%$ & 25 \\
    $\cos\theta_\mu\times p_\mu$ & $2\times2$ & $0.70/0.68$ & $6.1\%$ & $11.5\%$ & 69 \\
    $\cos\theta_\pi\times\theta_{\mu\pi}$ & $2\times3$ & $0.66/0.60$ & $5.7\%$ & $9.1\%$ & 4 \\
    $\cos\theta_\pi\times p_\pi$ & $2\times2$ & $0.53/0.51$ & $5.4\%$ & $17.2\%$ & 34 \\
    $\theta_p\times\delta p_T$ & $2\times2$ & $0.53/0.52$ & $8.4\%$ & $25.4\%$ & 63 \\
    $\theta_{\pi p}\times\delta p_T$ & $2\times2$ & $0.56/0.55$ & $10.7\%$ & $25.4\%$ & 58 \\
    $\theta_p\times W_{\pi p}$ & $2\times2$ & $0.40/0.39$ & $8.8\%$ & $39.4\%$ & 17 \\
    $\theta_{\pi p}\times W_{\pi p}$ & $2\times2$ & $0.21/0.20$ & $9.7\%$ & $39.4\%$ & 12 \\
    \bottomrule
  \end{tabular}
\end{table}

\subsection{Extractions}

All eight pairs are being extracted for FHC, RHC and the combined configuration, with generator
predictions produced as flattened histograms in the same bin order. At the time of writing the first
three are complete and already repeat the lesson of \S\ref{sec:c50}:
$\cos\theta_\pi\times\cos\theta_\mu$ gives $s_{\min}/s_1=0.245$ (FHC) and $0.315$ (RHC) with a
data-statistical term of $29$ and $25\%$, while $\cos\theta_\mu\times p_\mu$ --- the pair with the
best migration diagonals of the set --- comes out at $0.002$ in FHC, rank-deficient. A good diagonal
does not imply a usable regularised response. The completed table will be added when the array
finishes; the evaluation is \texttt{report/tools/extract\_eval.py}, which reads the multi-slice
sidecar where one exists.

% =============================================================================
\section{Reproduction}
\label{sec:repro}
% =============================================================================

\begin{small}
\begin{verbatim}
# screening (single pass, 1D and 2D; reproduces the published pre-screen exactly)
root -l -b -q 'macros/binning_screen_2d.C+("fhc","incl")'      # rhc; incl | 1p
root -l -b -q 'macros/tki_binning_scan.C("/data/uboone/processed/w",\
                     400,0.68,"thetap,thpipr",10)'
# theta_pipr without reprocessing (friend trees, for screening only)
logs/c50/run_pafriend.sh
# configurations
python3 scripts/make_bin_configs.py c50_candidates      # configs/c50
python3 scripts/make_bin_configs.py newobs_candidates   # configs/ + configs/d2
python3 scripts/check_ext_bins.py                       # generator vs configuration edges
# generator predictions (verifies the released histograms bit-for-bit)
generator_predictions/gen2d/regen_ext.sh
# reprocess with the opening angle, its gate, and promotion
./reprocess_thpipr.sh ; PROMOTE=1 ./reprocess_thpipr.sh
# extractions
sbatch --array=0-25%4 --export=ALL,REPO=$PWD,\
       MANIFEST=$PWD/slurm/c50_manifest.list    slurm/slurm_c50.sbatch
sbatch --array=0-29%6 --export=ALL,REPO=$PWD,\
       MANIFEST=$PWD/slurm/newobs_manifest.list slurm/slurm_extract.sbatch
# evaluation and figures
python3 report/tools/binning_eval_c50.py   # 0.50 candidates
python3 report/tools/extract_eval.py       # new observables and 2D
python3 report/tools/c50_figures.py
python3 report/tools/newobs_xsec_fig.py
\end{verbatim}
\end{small}

% =============================================================================
\section{Open items}
\label{sec:open}
% =============================================================================

\begin{enumerate}[itemsep=3pt]
  \item The two-dimensional extractions are running; their conditioning decides whether any pair is
    worth proposing as a released measurement. On the evidence so far only
    $\cos\theta_\pi\times\cos\theta_\mu$ is a candidate, and its data-statistical term is large.
  \item $\theta_p$ in RHC is rank-deficient. Either it is quoted for FHC and the combined
    configuration only, or the RHC binning is coarsened to three bins, which has not been screened.
  \item Two-bin $W_{\rm had}$ is better conditioned than the released six-bin configuration; a
    decision on replacing it is outstanding.
  \item The proton candidate is the leading true proton only $82\%$ of the time on signal events, and
    the truth-level observables use the leading proton. The effect is inside the response, but an
    inverted-proton-PID control region --- already a prerequisite for the signal region --- is what
    would constrain it.
  \item The released analysis note does not yet mention $\theta_p$ or $\theta_{\pi p}$; the binning
    appendix and the observable tables need entries once the decisions above are taken.
\end{enumerate}

\end{document}
